% Source: https://artofproblemsolving.com/community/c3h2861176_sub_hexagon_is_a_dependent_twitch_video_game_developed_by_terry_nacavagh_w
% Stable identity: aops-topic-2861176-q01
Each angle is $120^\circ$. Extend the sides $a_1$, $a_3$, $a_5$ until they meet. Each cut-off corner triangle has two angles of $180^\circ-120^\circ=60^\circ$, so it is equilateral, with side equal to the hexagon side it sits on ($a_2$, $a_4$ or $a_6$). The three extended lines form a large equilateral triangle.
\begin{asy}
usepackage("amsmath");
size(5cm);
defaultpen(fontsize(10pt));
real[] s={4,1,3,3,2,2};pair P=(0,0);pair[] V;
for(int i=0;i<6;++i){V.push(P);P=P+s[i]*dir(60*i);}
draw(V[0]--V[1]--V[2]--V[3]--V[4]--V[5]--cycle,linewidth(1));
pair Ta=extension(V[0],V[1],V[2],V[3]),Tb=extension(V[2],V[3],V[4],V[5]),Tc=extension(V[4],V[5],V[0],V[1]);
draw(V[1]--Ta--V[2],dashed);draw(V[3]--Tb--V[4],dashed);draw(V[5]--Tc--V[0],dashed);
for(int i=0;i<6;++i) label("$a_{"+string(i+1)+"}$",(V[i]+V[(i+1)%6])/2,-dir(60*i+90));
\end{asy}

Its sides are $a_6+a_1+a_2$, $a_2+a_3+a_4$ and $a_4+a_5+a_6$, and they are equal:
\[a_6+a_1+a_2=a_2+a_3+a_4=a_4+a_5+a_6.\]
The first equality gives $a_1-a_4=a_3-a_6$, and the second gives $a_3-a_6=a_5-a_2$. Taking absolute values, the positive differences of opposite sides are equal. $\blacksquare$
